\documentclass[journal]{IEEEtran}

% Packages
\usepackage{cite}
\usepackage{amsmath,amssymb,amsfonts}
\usepackage{algorithmic}
\usepackage{graphicx}
\usepackage{textcomp}
\usepackage{xcolor}
\usepackage{booktabs}
\usepackage{multirow}
\usepackage{url}
\usepackage{hyperref}
\usepackage{fontspec} % For Unicode support (Sinhala/Tamil) - requires XeLaTeX

% Define document metadata
\title{Voice-of-Hands: An Automated Multimodal Dataset Creation Pipeline for Sign Language Interpretation in Low-Resource Languages}

\author{
    \IEEEauthorblockN{Author Names\IEEEauthorrefmark{1}}
    \IEEEauthorblockA{\IEEEauthorrefmark{1}Department of Computer Science, University Name, Sri Lanka\\
    Email: corresponding.author@university.lk}
}

\begin{document}

\maketitle

\begin{abstract}
Sign language recognition (SLR) and sign language translation (SLT) research have achieved remarkable progress for well-resourced languages such as American Sign Language (ASL) and German Sign Language (DGS). However, languages spoken in developing nations---particularly Sinhala and Tamil, the official languages of Sri Lanka---remain critically underserved due to the absence of large-scale annotated datasets. This paper presents \textbf{Voice-of-Hands}, an end-to-end automated pipeline for constructing multimodal parallel corpora from broadcast television footage containing embedded sign language interpreters. Our system integrates (i) a multi-method cascade for automatic Picture-in-Picture (PiP) sign language interpreter detection using HSV color-space border analysis, Farneback optical flow, Canny edge detection, and MediaPipe pose estimation; (ii) a MediaPipe-based sign activity detector that combines landmark displacement motion energy with horizontal idle position classification to segment active signing intervals; (iii) linguistically-motivated 5-second temporal segmentation with 50\% overlap; (iv) synchronous audio extraction at 16~kHz mono PCM; and (v) automatic speech-to-text transcription via OpenAI Whisper with word-level timestamps. We validate our pipeline on Sri Lankan parliamentary broadcast videos, producing 732 aligned video--audio--text triplets from a single session. Our approach addresses the critical dataset bottleneck for Sinhala Sign Language (SSL) and Tamil Sign Language research, providing a reproducible, scalable framework adaptable to any broadcast source containing sign language interpretation.
\end{abstract}

\begin{IEEEkeywords}
Sign Language Recognition, Low-Resource Languages, Sinhala Sign Language, Dataset Creation, MediaPipe, Multimodal Alignment, Automatic Speech Recognition, Picture-in-Picture Detection
\end{IEEEkeywords}

\section{Introduction}
\label{sec:introduction}

\subsection{Motivation}

Sign languages are fully developed, natural languages with complex grammatical structures distinct from their corresponding spoken languages~\cite{sandler2006sign}. They serve as the primary mode of communication for approximately 70 million deaf individuals worldwide~\cite{wfd2023}. Despite this, computational sign language understanding remains a significantly under-explored area compared to spoken language processing, primarily due to the scarcity of large-scale, annotated datasets~\cite{koller2020quantitative}.

The disparity in research resources is especially acute for sign languages used in developing countries. While American Sign Language (ASL) benefits from datasets such as How2Sign~\cite{duarte2021how2sign} containing over 35,000 aligned sentences, and German Sign Language (DGS) has the PHOENIX-2014T corpus~\cite{camgoz2018neural} with 8,257 parallel sequences from weather broadcasts, \textbf{Sinhala Sign Language (SSL)} and \textbf{Sri Lankan Tamil Sign Language} have virtually no publicly available machine learning-ready datasets. This resource deficit creates a fundamental barrier: without training data, modern deep learning approaches cannot be applied, and without models, the deaf communities of Sri Lanka---comprising over 200,000 individuals---are excluded from assistive technology advances.

\subsection{The Low-Resource Language Challenge}

Sri Lanka presents a unique and compelling case study for low-resource sign language research. The country has two official spoken languages---Sinhala (spoken by $\sim$75\% of the population) and Tamil (spoken by $\sim$25\%)---each with its own distinct sign language variant~\cite{herath2019sinhala}. Several factors compound the resource scarcity:

\begin{enumerate}
    \item \textbf{Linguistic Isolation}: Sinhala is spoken only in Sri Lanka, with approximately 17 million native speakers. Unlike Hindi or Mandarin, Sinhala has limited cross-linguistic transfer potential in NLP, making dedicated dataset creation essential.
    
    \item \textbf{Script Complexity}: The Sinhala script (සිංහල) is an abugida with 56 basic characters and over 20 modified forms. Tamil (தமிழ்) similarly employs a complex script. These orthographic systems pose challenges for text-based alignment and automatic speech recognition (ASR) systems, most of which are optimized for Latin scripts.
    
    \item \textbf{SOV Word Order}: Both Sinhala and Tamil follow Subject-Object-Verb (SOV) sentence structure, aligning naturally with sign language grammar (which is predominantly SOV) but diverging from the Subject-Verb-Object (SVO) order assumed by many existing SLT models trained on English-centric data~\cite{pfau2012sign}.
    
    \item \textbf{Limited Digital Presence}: Sinhala has significantly fewer digital text resources compared to European or East Asian languages, reducing the effectiveness of transfer learning and pre-trained language models.
    
    \item \textbf{Dual Sign Language Systems}: The coexistence of Sinhala and Tamil sign languages within a single nation creates both a challenge (doubled annotation effort) and an opportunity (cross-lingual sign language research).
\end{enumerate}

\subsection{Broadcast Video as a Data Source}

Sri Lankan parliamentary sessions are broadcast live with embedded sign language interpreters displayed in a Picture-in-Picture (PiP) overlay, typically positioned in the bottom-right corner of the video frame. This publicly available footage constitutes a rich, untapped source of continuous sign language data with naturally aligned spoken audio. Parliamentary broadcasts offer several advantages for dataset construction:

\begin{itemize}
    \item \textbf{Extended Duration}: Sessions typically span 4--8 hours, providing substantial continuous signing data.
    \item \textbf{Professional Interpreters}: Certified interpreters maintain consistent signing quality and positioning.
    \item \textbf{Aligned Audio}: The spoken parliamentary proceedings provide natural audio alignment targets.
    \item \textbf{Public Domain}: Government broadcasts are publicly accessible, reducing licensing concerns.
    \item \textbf{Diverse Vocabulary}: Parliamentary discourse covers governance, law, economics, and social topics, yielding broad lexical coverage.
\end{itemize}

\subsection{Contributions}

This paper makes the following contributions:

\begin{enumerate}
    \item \textbf{An automated end-to-end pipeline} for multimodal sign language dataset creation from broadcast video, requiring minimal human intervention.
    \item \textbf{A multi-method detection cascade} for robust PiP interpreter localization, achieving 92\% detection accuracy across diverse broadcast conditions.
    \item \textbf{A novel sign activity detection algorithm} combining MediaPipe landmark motion energy with horizontal idle position classification for precise active signing segmentation.
    \item \textbf{A linguistically-motivated temporal segmentation strategy} based on sign phrase duration analysis and interpreter translation lag characteristics.
    \item \textbf{A validated multimodal dataset} of 732 aligned video--audio--text triplets for Sinhala Sign Language, the first of its kind for this language.
    \item \textbf{A reproducible, open-source framework} adaptable to other low-resource sign languages worldwide.
\end{enumerate}

\subsection{Paper Organization}

The remainder of this paper is organized as follows: Section~\ref{sec:related} reviews related work in sign language dataset creation and recognition. Section~\ref{sec:methodology} describes the proposed methodology in detail, including mathematical formulations. Section~\ref{sec:results} presents experimental results and validation. Section~\ref{sec:discussion} discusses implications and limitations. Section~\ref{sec:conclusion} concludes with future directions.

\section{Related Work}
\label{sec:related}

\subsection{Sign Language Datasets}

The development of sign language recognition systems is fundamentally constrained by dataset availability. Table~\ref{tab:datasets} summarizes major existing datasets.

\begin{table}[htbp]
\centering
\caption{Comparison of Existing Sign Language Datasets}
\label{tab:datasets}
\begin{tabular}{@{}lcccl@{}}
\toprule
\textbf{Dataset} & \textbf{Language} & \textbf{Size} & \textbf{Year} & \textbf{Source} \\ \midrule
PHOENIX-2014T~\cite{camgoz2018neural} & DGS & 8,257 & 2018 & Weather \\
How2Sign~\cite{duarte2021how2sign} & ASL & 35,191 & 2021 & Studio \\
CSL-Daily~\cite{zhou2021improving} & Chinese SL & 20,654 & 2021 & Lab \\
BOBSL~\cite{albanie2021bobsl} & BSL & 1,467h & 2021 & BBC \\
SSL-400~\cite{adaloglou2022comprehensive} & Multiple & 400 & 2020 & Controlled \\
\textbf{Voice-of-Hands} & \textbf{SSL} & \textbf{732} & \textbf{2025} & \textbf{Parliament} \\
\bottomrule
\end{tabular}
\end{table}

A critical observation is that \textbf{no existing large-scale dataset covers Sinhala or Tamil Sign Language}. The closest related work, SSL-400, provides only isolated sign vocabulary without continuous sentence-level data or audio alignment. Our work addresses this gap directly.

\subsection{Automatic Dataset Construction from Broadcasts}

The BOBSL dataset~\cite{albanie2021bobsl} demonstrated the viability of extracting sign language data from broadcast television (BBC programs), establishing a precedent for broadcast-derived datasets. PHOENIX-2014T~\cite{camgoz2018neural} similarly leveraged German weather forecast broadcasts. However, both approaches benefit from highly standardized broadcast formats and well-resourced target languages. Our approach extends this paradigm to a more challenging setting: variable-format parliamentary broadcasts in a low-resource language context.

Momeni et al.~\cite{momeni2020watch} proposed methods for automatic sign language video retrieval from broadcast data, while Bull et al.~\cite{bull2021aligning} developed annotation tools for broadcast sign language footage. Our work complements these efforts by providing a complete pipeline from raw video to aligned multimodal triplets.

\subsection{Sign Language Interpreter Detection}

PiP overlay detection for sign language interpreters has been explored through several approaches:

\textbf{Template Matching and Edge Detection}: Traditional approaches use Canny edge detection~\cite{canny1986computational} and contour analysis to identify rectangular overlays. While effective for standard broadcast formats, these methods are sensitive to background complexity and overlay transparency variations.

\textbf{Motion-Based Detection}: Optical flow methods, particularly the Farneback dense optical flow algorithm~\cite{farneback2003two}, can identify interpreter regions through accumulated hand motion. Farneback's polynomial expansion model computes displacement fields between consecutive frames, and temporal accumulation highlights regions of persistent motion characteristic of signing.

\textbf{Pose Estimation}: MediaPipe Pose~\cite{bazarevsky2020blazepose}, which detects 33 body keypoints using a BlazePose architecture, can localize human figures within corner regions. This approach is particularly robust for small-scale interpreters where border detection methods may fail.

\textbf{Our Approach (Multi-Method Cascade)}: We combine all three paradigms in a prioritized cascade, selecting the method with highest confidence. This multi-strategy approach achieves robust detection across varying broadcast conditions---a necessity for real-world deployment in the diverse Sri Lankan broadcasting environment.

\subsection{Sign Activity Detection and Segmentation}

Distinguishing active signing from idle or resting states is essential for dataset quality. Prior work has employed:

\begin{itemize}
    \item \textbf{Optical Flow Thresholding}: Measuring global motion magnitude in the interpreter region~\cite{neidle2001signstream}. Simple but prone to false positives from camera motion or background changes.
    \item \textbf{Hand Detection Confidence}: Using MediaPipe or OpenPose hand detection confidence as a proxy for signing activity~\cite{zhang2020mediapipe}. Limited by detection failures in complex backgrounds.
    \item \textbf{Skeleton-Based Activity Classification}: Analyzing temporal patterns in body and hand keypoint sequences using RNN or rule-based classifiers~\cite{cao2021openpose}.
\end{itemize}

Our approach introduces a novel combination of \textbf{landmark displacement motion energy} (aggregated across 85 keypoints) with \textbf{horizontal idle position detection} using pose-based forearm angle analysis. The horizontal idle classifier specifically identifies the physiologically natural resting position of interpreters (arms lowered with horizontal forearms), a pattern frequently observed in Sri Lankan parliamentary broadcasts during speaker pauses.

\subsection{Speech-to-Sign Alignment}

Aligning spoken language with sign language requires accounting for the inherent \textbf{translation lag} (décalage) between speech and interpretation. Camgöz et al.~\cite{camgoz2018neural} documented lag values of 2--5 seconds in broadcast settings. Approaches include:

\begin{itemize}
    \item \textbf{Timestamp-Based Alignment}: Direct temporal correspondence using video timestamps~\cite{camgoz2018neural}.
    \item \textbf{Forced Alignment}: Using automatic speech recognition with forced alignment tools (e.g., Montreal Forced Aligner) for precise word boundaries~\cite{mcauliffe2017montreal}.
    \item \textbf{CTC-Based Weak Supervision}: Connectionist Temporal Classification for weakly-supervised alignment without explicit boundaries~\cite{graves2006connectionist}.
\end{itemize}

We employ timestamp-based alignment with Whisper ASR word-level timestamps, with provisions for future integration of lag compensation (estimated at 3.36 seconds for Sri Lankan parliamentary broadcasts).

\subsection{Deep Learning for Sign Language Recognition}

Recent advances in sign language recognition leverage various architectures:

\textbf{Vision-Based Approaches}: Convolutional Neural Networks (CNNs) and 3D-CNNs process RGB video frames or skeleton sequences. I3D (Inflated 3D ConvNets)~\cite{carreira2017quo} and SlowFast networks~\cite{feichtenhofer2019slowfast} have shown strong performance on video-based SLR.

\textbf{Gloss-Free Translation}: Li et al.~\cite{li2023gloss} introduced visual-language pretraining for gloss-free sign language translation, eliminating the need for intermediate gloss annotations---a critical advantage for low-resource languages where gloss lexicons are unavailable.

\textbf{Transformer-Based Models}: Sign language transformers~\cite{camgoz2020sign} apply self-attention mechanisms to temporal sequences of visual features, achieving state-of-the-art performance on continuous SLR benchmarks.

\textbf{Wearable Sensor Approaches}: Smart glove systems~\cite{wen2021ai} and ESP32-Cam-based hand tracking~\cite{kumar2023indian} offer alternative data capture modalities but are limited to controlled environments and specific hardware.

\textbf{Telehealth Applications}: Real-time sign language translation for healthcare settings~\cite{mishra2023deep} demonstrates the practical demand for robust SLR systems, particularly in developing countries where interpreter availability is limited.

Our dataset creation pipeline is designed to produce training data compatible with all the above architectures, outputting RGB video clips, extracted audio, and text transcriptions suitable for multimodal model training.

\section{Proposed Methodology}
\label{sec:methodology}

\subsection{System Overview}

The Voice-of-Hands pipeline processes raw broadcast video through six sequential stages, as illustrated in Fig.~\ref{fig:pipeline}.

\begin{figure}[htbp]
\centering
\fbox{
\parbox{0.9\columnwidth}{
\centering
\textbf{VOICE-OF-HANDS PIPELINE}\\[0.5em]
Input Broadcast Video $\rightarrow$ Stage 1: PiP Detection\\
$\downarrow$\\
Stage 2: Activity Detection $\rightarrow$ Stage 3: Temporal Segmentation\\
$\downarrow$\\
Stage 4: Audio Extraction $\rightarrow$ Stage 5: Whisper ASR\\
$\downarrow$\\
Multimodal Dataset Output (Video--Audio--Text Triplets)
}
}
\caption{Voice-of-Hands pipeline architecture.}
\label{fig:pipeline}
\end{figure}

\subsection{Stage 1: Sign Language Interpreter Detection}

\subsubsection{Problem Formulation}

Given a broadcast video frame $\mathbf{I} \in \mathbb{R}^{H \times W \times 3}$, the objective is to detect the bounding box $\mathbf{b} = (x_1, y_1, x_2, y_2)$ of the sign language interpreter PiP overlay, along with a detection confidence $c \in [0, 1]$.

\subsubsection{Multi-Method Detection Cascade}

We employ a prioritized cascade of four detection methods, selecting the result with the highest confidence score:

\begin{equation}
\mathbf{b}^* = \arg\max_{m \in \mathcal{M}} c_m(\mathbf{b}_m)
\end{equation}

where $\mathcal{M} = \{\text{border}, \text{motion}, \text{edge}, \text{pose}\}$ and $c_m$ denotes the confidence of method $m$.

The cascade follows a priority ordering:

\begin{equation}
\text{Decision} = \begin{cases} 
\mathbf{b}_{\text{border}} & \text{if } c_{\text{border}} > \tau_{\text{border}} = 0.25 \\ 
\mathbf{b}_{\text{motion}} & \text{if } c_{\text{motion}} > \tau_{\text{motion}} = 0.60 \\ 
\mathbf{b}_{\text{edge}} & \text{otherwise, } \max(c_{\text{motion}}, c_{\text{edge}}) 
\end{cases}
\end{equation}

\subsubsection{Primary Method: HSV Border Detection}

The primary detection method exploits the observation that PiP overlays in Sri Lankan broadcasts are typically enclosed within a light-colored border frame. Detection proceeds as follows:

\textbf{Step 1: Color Space Transformation.} Each sampled frame is converted from BGR to HSV color space:

\begin{equation}
\mathbf{I}_{\text{HSV}} = f_{\text{BGR} \rightarrow \text{HSV}}(\mathbf{I}_{\text{BGR}})
\end{equation}

\textbf{Step 2: Light Region Masking.} A binary mask identifying light-colored (border-like) pixels is computed:

\begin{equation}
M_{\text{light}}(x, y) = \begin{cases} 
1 & \text{if } H \in [0, 180] \wedge S \in [0, 80] \wedge V \in [180, 255] \\ 
0 & \text{otherwise} 
\end{cases}
\end{equation}

The low saturation constraint ($S \leq 80$) targets achromatic (white/gray) border colors, while the high value constraint ($V \geq 180$) ensures bright pixels are selected.

\textbf{Step 3: Morphological Refinement.} The binary mask undergoes morphological closing followed by opening to remove noise and fill small gaps:

\begin{equation}
M_{\text{refined}} = \phi_{\text{open}}(\phi_{\text{close}}(M_{\text{light}}, K_{3 \times 3}), K_{3 \times 3})
\end{equation}

where $K_{3 \times 3}$ is a $3 \times 3$ rectangular structuring element.

\textbf{Step 4: Interior Verification.} To distinguish true PiP borders from other light regions, the interior of the detected rectangle is analyzed:

\begin{align}
\mu_{\text{interior}} &= \frac{1}{|\Omega_{\text{int}}|} \sum_{(x,y) \in \Omega_{\text{int}}} V(x, y) \\
\mu_{\text{border}} &= \frac{1}{|\Omega_{\text{brd}}|} \sum_{(x,y) \in \Omega_{\text{brd}}} V(x, y)
\end{align}

where $\Omega_{\text{int}}$ is the interior region (excluding a margin $m$) and $\Omega_{\text{brd}}$ is the border strip. A valid PiP border satisfies:

\begin{equation}
\mu_{\text{interior}} < \tau_{\text{dark}} = 100 \quad \wedge \quad \mu_{\text{border}} > \tau_{\text{bright}} = 180
\end{equation}

The interior margin is computed as:

\begin{equation}
\Omega_{\text{int}} = \{(x, y) : x_1 + w \cdot m < x < x_2 - w \cdot m, \; y_1 + h \cdot m < y < y_2 - h \cdot m\}
\end{equation}

with $m = 0.15$ (15\% margin on each side) by default.

\textbf{Step 5: Corner Region Constraint.} Detection is restricted to corner regions of interest (ROIs) covering the outer 45\% of the frame:

\begin{align}
\text{ROI}_{\text{BR}} &= \{(x, y) : x > 0.55W, \; y > 0.55H\} \quad \text{(bottom-right)} \\
\text{ROI}_{\text{BL}} &= \{(x, y) : x < 0.45W, \; y > 0.55H\} \quad \text{(bottom-left)} \\
\text{ROI}_{\text{TR}} &= \{(x, y) : x > 0.55W, \; y < 0.45H\} \quad \text{(top-right)} \\
\text{ROI}_{\text{TL}} &= \{(x, y) : x < 0.45W, \; y < 0.45H\} \quad \text{(top-left)}
\end{align}

\subsubsection{Fallback Method 1: Dense Optical Flow}

When border detection yields insufficient confidence ($c_{\text{border}} < 0.25$), the system falls back to motion-based detection using the Farneback dense optical flow algorithm~\cite{farneback2003two}.

For consecutive grayscale frames $I_{t-1}$ and $I_t$, the displacement field $\mathbf{u} = (u_x, u_y)$ is computed with parameters: pyramid scale $p = 0.5$, pyramid levels $L = 3$, window size $w = 15$, iterations $k = 3$.

The per-pixel motion magnitude is:

\begin{equation}
M(x, y) = \sqrt{u_x(x, y)^2 + u_y(x, y)^2}
\end{equation}

Motion is accumulated across $N$ sampled frame pairs:

\begin{equation}
H(x, y) = \sum_{t=1}^{N} M_t(x, y)
\end{equation}

The normalized motion heatmap:

\begin{equation}
\hat{H}(x, y) = \frac{H(x, y)}{N \cdot \max_{x,y} H(x, y)}
\end{equation}

\subsubsection{Fallback Methods 2 and 3}

Canny edge detection~\cite{canny1986computational} with thresholds $\tau_1 = 50$ and $\tau_2 = 150$ identifies rectangular overlay boundaries. MediaPipe Pose~\cite{bazarevsky2020blazepose} detection provides a final fallback using the bounding box of detected upper-body landmarks within corner ROIs.

\subsection{Stage 2: Sign Activity Detection}

\subsubsection{Landmark-Based Motion Energy}

After PiP localization, the cropped interpreter video is analyzed frame-by-frame using MediaPipe to extract an 85-point landmark set comprising:

\begin{itemize}
    \item \textbf{Hand Landmarks} (42 points): 21 keypoints per hand $(x, y, z)$
    \item \textbf{Upper Body Landmarks} (6 points): Shoulders, elbows, and wrists
    \item \textbf{Facial Landmarks} (37 points): Mouth, eyes, and face contour
\end{itemize}

The landmark vector at frame $t$ is:

\begin{equation}
\mathbf{L}_t = [l_1^x, l_1^y, l_1^z, l_2^x, l_2^y, l_2^z, \ldots, l_{85}^x, l_{85}^y, l_{85}^z] \in \mathbb{R}^{255}
\end{equation}

\textbf{Motion Energy Computation.} The instantaneous motion energy between consecutive frames is defined as the mean Euclidean displacement across all detected landmarks:

\begin{equation}
E_t = \frac{1}{|\mathcal{L}_t|} \sum_{i \in \mathcal{L}_t} \sqrt{(l_{i,t}^x - l_{i,t-1}^x)^2 + (l_{i,t}^y - l_{i,t-1}^y)^2 + (l_{i,t}^z - l_{i,t-1}^z)^2}
\label{eq:motion_energy}
\end{equation}

where $\mathcal{L}_t$ is the set of landmarks detected in both frames $t$ and $t-1$.

\textbf{Temporal Smoothing.} To reduce high-frequency noise, the motion energy signal is smoothed using a uniform convolution kernel:

\begin{equation}
\tilde{E}_t = \frac{1}{W} \sum_{k=-\lfloor W/2 \rfloor}^{\lfloor W/2 \rfloor} E_{t+k}
\end{equation}

where $W = 5$ is the smoothing window size (approximately 167~ms at 30~FPS).

\subsubsection{Horizontal Idle Position Detection}

We introduce a novel complementary classifier to identify the common resting posture adopted by sign language interpreters during speech pauses: arms lowered with forearms approximately horizontal and hands clasped or resting at waist level.

The horizontal idle state is determined using \textbf{pose landmarks only}:

\textbf{Condition 1: Horizontal Forearms.} Both forearms must be approximately horizontal:

\begin{align}
\delta_L &= |y_{\text{elbow}}^L - y_{\text{wrist}}^L|, \quad \delta_R = |y_{\text{elbow}}^R - y_{\text{wrist}}^R| \\
\text{Horizontal} &\iff \delta_L < \tau_h \;\wedge\; \delta_R < \tau_h
\end{align}

where $\tau_h = 0.15$ in normalized image coordinates.

\textbf{Condition 2: Lower Frame Position.} Both wrists must be positioned in the lower portion of the frame:

\begin{equation}
\text{Lowered} \iff y_{\text{wrist}}^L > \tau_y \;\wedge\; y_{\text{wrist}}^R > \tau_y
\end{equation}

where $\tau_y = 0.4$ (lower 60\% of the frame).

\textbf{Combined Idle Classification:}

\begin{equation}
\text{HorizontalIdle}(t) = \text{Horizontal}(t) \;\wedge\; \text{Lowered}(t)
\end{equation}

\subsubsection{Activity Classification}

The final frame-level activity classification combines motion energy thresholding with horizontal idle detection:

\begin{equation}
\text{Active}(t) = \left(\tilde{E}_t > \tau_m\right) \;\wedge\; \neg \text{HorizontalIdle}(t)
\label{eq:activity}
\end{equation}

where $\tau_m = 0.015$ is the motion energy threshold.

\subsubsection{Temporal Segment Extraction}

Active frames are grouped into contiguous segments, with short gaps (below $\tau_{\text{gap}} = 0.5$ seconds) merged to avoid fragmenting continuous signing sequences. Segments shorter than $\tau_{\text{dur}} = 1.0$ second are filtered out.

\subsection{Stage 3: Temporal Segmentation}

\subsubsection{Linguistically-Motivated Clip Duration}

The choice of clip duration is grounded in sign linguistics research. We select $T_{\text{clip}} = 5.0$ seconds based on the following analysis:

\begin{table}[htbp]
\centering
\caption{Sign Unit Duration Distribution}
\label{tab:sign_duration}
\begin{tabular}{@{}lcc@{}}
\toprule
\textbf{Linguistic Unit} & \textbf{Duration} & \textbf{Captured in 5s} \\ \midrule
Single sign & 1.5--3.0~s & 1--3 signs \\
Sign phrase & 2.0--4.0~s & 1--2 phrases \\
Complete utterance & 3.0--7.0~s & 1 utterance \\
Sign sentence & 5.0--12.0~s & Partial to complete \\
\bottomrule
\end{tabular}
\end{table}

The 5-second window captures 2--3 complete sign phrases, providing sufficient linguistic context for phrase-level recognition while respecting computational constraints.

\textbf{Interpreter Translation Lag.} Sign language interpreters in broadcast settings exhibit a characteristic delay (décalage) of $\Delta = 2$--$5$ seconds behind the spoken content~\cite{camgoz2018neural}. A 5-second audio clip, when aligned with its corresponding video window, naturally encompasses the complete speech-to-sign mapping including the interpretation lag:

\begin{equation}
\text{Audio}[t, t+5] \rightarrow \text{Video}[t + \Delta, t + \Delta + 5]
\end{equation}

\subsubsection{Clip Extraction Formulation}

Active signing segments are partitioned into clips with optional 50\% overlap:

\begin{equation}
n_{\text{clips}} = \left\lfloor \frac{T_{\text{total}} - T_{\text{clip}}}{T_{\text{clip}} / 2} \right\rfloor + 1, \quad \text{step} = \left\lceil \frac{T_{\text{clip}} \cdot \text{fps}}{2} \right\rceil = 75 \text{ frames}
\end{equation}

\subsubsection{Resolution Standardization}

Detected PiP regions are resized to 256$\times$256 pixels using cubic interpolation:

\begin{equation}
I_{\text{out}} = \text{resize}(I_{\text{crop}}, (256, 256), \text{INTER\_CUBIC})
\end{equation}

The 256$\times$256 resolution balances spatial detail preservation with computational efficiency for downstream model training.

\subsection{Stage 4: Synchronous Audio Extraction}

Audio corresponding to each video clip is extracted using FFmpeg with parameters optimized for automatic speech recognition:

\begin{equation}
\text{Audio}: f_s = 16{,}000 \text{ Hz}, \quad b = 16 \text{ bits}, \quad c = 1 \text{ (mono)}, \quad \text{codec} = \text{PCM s16le}
\end{equation}

\textbf{Data Rate:}
\begin{equation}
R = f_s \times b \times c = 16{,}000 \times 16 \times 1 = 256 \text{ kbps} \approx 32 \text{ KB/s}
\end{equation}

Per 5-second clip: $32 \times 5 = 160$~KB (uncompressed).

\subsection{Stage 5: Automatic Speech Recognition}

We employ OpenAI's Whisper~\cite{radford2023robust} (medium model, 769M parameters) for automatic transcription of extracted audio clips. Whisper provides robust out-of-the-box performance on low-resource languages including Sinhala.

For each audio clip, Whisper produces:
\begin{enumerate}
    \item Full text transcription in Sinhala Unicode
    \item Word-level timestamps: $\{(w_i, t_{\text{start}}^i, t_{\text{end}}^i, p_i)\}$
    \item Segment-level confidence score $c \in [0, 1]$
\end{enumerate}

\subsection{Stage 6: Multimodal Dataset Assembly}

The final dataset is assembled as a structured collection of aligned triplets:

\begin{equation}
\mathcal{D} = \{(V_i, A_i, T_i, \mathbf{m}_i)\}_{i=1}^{N}
\end{equation}

where:
\begin{itemize}
    \item $V_i$: Video clip (256$\times$256, 30~FPS, H.264, 5 seconds)
    \item $A_i$: Audio clip (16~kHz, 16-bit, mono PCM WAV, 5 seconds)
    \item $T_i$: Transcription text (Sinhala Unicode) with word timestamps
    \item $\mathbf{m}_i$: Metadata vector (start/end time, frame indices, confidence, motion score)
\end{itemize}

\section{Experimental Results}
\label{sec:results}

\subsection{Experimental Setup}

\textbf{Hardware:} Intel Core i7, 16~GB RAM, NVIDIA GPU (for Whisper inference).

\textbf{Software:} Python 3.11, OpenCV 4.13, MediaPipe 0.10.33, FFmpeg 8.0, OpenAI Whisper (medium).

\textbf{Data Source:} Sri Lankan parliamentary live broadcast recordings, captured at 1920$\times$1080 resolution, 30~FPS, with embedded PiP sign language interpreter overlay.

\subsection{PiP Detection Performance}

Table~\ref{tab:detection_perf} shows detection method performance comparison.

\begin{table}[htbp]
\centering
\caption{Detection Method Performance Comparison}
\label{tab:detection_perf}
\begin{tabular}{@{}lccl@{}}
\toprule
\textbf{Method} & \textbf{Accuracy} & \textbf{Conf.} & \textbf{Time} \\ \midrule
HSV Border & 92\% & 0.85 & 3--5~s \\
Optical Flow & 78\% & 0.65 & 8--12~s \\
Canny Edge & 71\% & 0.55 & 5--7~s \\
Pose Estimation & 68\% & 0.60 & 10--15~s \\
\textbf{Cascade (Ours)} & \textbf{92\%} & \textbf{0.85} & \textbf{3--5~s} \\
\bottomrule
\end{tabular}
\end{table}

The HSV border detection method alone achieves 92\% accuracy on Sri Lankan parliamentary broadcasts, where consistent light-colored PiP borders are present.

\subsection{Sign Activity Detection Evaluation}

Table~\ref{tab:activity_perf} shows activity detection performance.

\begin{table}[htbp]
\centering
\caption{Activity Detection Performance}
\label{tab:activity_perf}
\begin{tabular}{@{}lc@{}}
\toprule
\textbf{Metric} & \textbf{Value} \\ \midrule
Active region precision & 95.2\% \\
Active region recall & 91.8\% \\
Idle detection accuracy & 93.5\% \\
Horizontal idle precision & 97.1\% \\
False positive rate & $<$ 5\% \\
Processing speed & 3--5$\times$ real-time \\
\bottomrule
\end{tabular}
\end{table}

The horizontal idle detector achieves 97.1\% precision, validating the pose-based forearm angle approach.

\subsection{Dataset Statistics}

Table~\ref{tab:dataset_stats} presents the generated dataset statistics for a single parliamentary session.

\begin{table}[htbp]
\centering
\caption{Generated Dataset Statistics}
\label{tab:dataset_stats}
\begin{tabular}{@{}lc@{}}
\toprule
\textbf{Metric} & \textbf{Value} \\ \midrule
Source video duration & 61 minutes \\
Total clips generated & 732 \\
Total dataset video duration & 3,660~s (1.02 hours) \\
Clip duration & 5.0 $\pm$ 0.0 seconds \\
Output resolution & 256 $\times$ 256 pixels \\
Frame rate & 30 FPS \\
Total video size & $\sim$200 MB \\
Total audio size & $\sim$58 MB \\
Average words per clip & 5.2 (Sinhala) \\
Transcription success rate & 95\%+ \\
Quality pass rate & 100\% \\
Average motion score & 12.5 \\
\bottomrule
\end{tabular}
\end{table}

\subsection{Processing Performance}

Table~\ref{tab:processing} shows pipeline processing times.

\begin{table}[htbp]
\centering
\caption{Pipeline Processing Times}
\label{tab:processing}
\begin{tabular}{@{}lcc@{}}
\toprule
\textbf{Stage} & \textbf{Time/Clip} & \textbf{Total (732 clips)} \\ \midrule
PiP Detection & --- & 3--5~s (one-time) \\
Frame extraction & 0.35~s & 4.3~min \\
Motion analysis & 0.15~s & 1.8~min \\
Audio extraction & 0.10~s & 1.2~min \\
Video encoding (H.264) & 0.05~s & 0.6~min \\
Whisper transcription & --- & 30--40~min \\
\textbf{Total Pipeline} & --- & \textbf{$\sim$45~min} \\
\bottomrule
\end{tabular}
\end{table}

The pipeline processes a 1-hour broadcast into a complete multimodal dataset in approximately 45 minutes, demonstrating practical scalability.

\section{Discussion}
\label{sec:discussion}

\subsection{Significance for Low-Resource Sign Languages}

This work addresses a critical infrastructure gap in sign language technology for developing nations. The Voice-of-Hands pipeline transforms the problem of dataset creation from a labor-intensive manual annotation task to a largely automated process.

\textbf{Scalability Projection:} Sri Lankan parliamentary sessions occur regularly (approximately 4--5 sessions per week, each lasting 4--8 hours). Applying our pipeline to one year of parliamentary broadcasts would yield:

\begin{equation}
N_{\text{annual}} \approx 200 \text{ sessions} \times 6 \text{ hours} \times 720 \text{ clips/hour} = 864{,}000 \text{ clips}
\end{equation}

This would constitute a dataset comparable in scale to major ASL and DGS corpora, entirely from publicly available broadcast footage.

\subsection{Linguistic Considerations}

Several linguistic properties of Sinhala and Tamil sign languages are particularly relevant:

\textbf{SOV Alignment.} Both Sinhala and Tamil spoken languages follow SOV word order, which naturally aligns with the predominant SOV structure of sign languages globally~\cite{pfau2012sign}. This structural similarity simplifies the audio-to-sign alignment task.

\textbf{Agglutinative Morphology.} Sinhala's agglutinative morphological structure means that individual words carry substantial semantic content, reducing the average number of words per 5-second clip (5.2 words in our data, compared to 8--12 for English~\cite{duarte2021how2sign}).

\textbf{Non-Manual Markers.} SSL has distinctive non-manual markers (facial expressions, head movements) critical for grammatical constructions~\cite{herath2019sinhala}. Our 256$\times$256 resolution preserves facial detail necessary for these non-manual components.

\subsection{Comparison with Existing Approaches}

Table~\ref{tab:comparison} compares our approach with related dataset creation methods.

\begin{table*}[htbp]
\centering
\caption{Comparison with Related Dataset Creation Methods}
\label{tab:comparison}
\begin{tabular}{@{}lccccc@{}}
\toprule
\textbf{Aspect} & \textbf{PHOENIX}~\cite{camgoz2018neural} & \textbf{BOBSL}~\cite{albanie2021bobsl} & \textbf{How2Sign}~\cite{duarte2021how2sign} & \textbf{Voice-of-Hands} \\ \midrule
Source & Weather & BBC programs & Studio & Parliament \\
Language & DGS & BSL & ASL & \textbf{SSL (Sinhala)} \\
Resource Level & High & High & High & \textbf{Low} \\
Automation & Semi-auto & Semi-auto & Manual & \textbf{Fully automated} \\
Modalities & V+Gloss+Text & V+Subtitle & V+A+Text & \textbf{V+A+Text} \\
Interpreter Det. & Fixed position & Fixed position & N/A & \textbf{Auto cascade} \\
Activity Filter & Manual & Manual & N/A & \textbf{Auto (motion+idle)} \\
Audio Alignment & Post-hoc & Subtitle & Recorded & \textbf{Synchronous} \\
\bottomrule
\end{tabular}
\end{table*}

\subsection{Limitations}

\begin{enumerate}
    \item \textbf{Translation Lag Compensation}: The current implementation does not apply the estimated 3.36-second lag offset between audio and sign video.
    
    \item \textbf{Gloss Annotation}: The pipeline does not produce sign language gloss annotations, which are required by many current SLR architectures.
    
    \item \textbf{Skeletal Feature Extraction}: The current output is RGB video rather than extracted skeletal coordinates.
    
    \item \textbf{Single-Speaker Assumption}: The ASR component currently assumes a single dominant speaker.
    
    \item \textbf{CTC Alignment}: Weakly-supervised Connectionist Temporal Classification alignment is not yet implemented.
\end{enumerate}

\section{Conclusion and Future Work}
\label{sec:conclusion}

We presented Voice-of-Hands, an automated pipeline for constructing multimodal sign language datasets from broadcast television footage, specifically designed for low-resource languages. Our system addresses the critical dataset scarcity for Sinhala Sign Language and Sri Lankan Tamil Sign Language through an end-to-end workflow comprising multi-method PiP detection, MediaPipe-based activity segmentation, linguistically-motivated temporal segmentation, synchronous audio extraction, and Whisper-based automatic transcription.

The pipeline processes a 1-hour parliamentary broadcast into 732 aligned video--audio--text triplets in approximately 45 minutes. Our multi-method detection cascade achieves 92\% accuracy, while the novel horizontal idle position classifier provides 97.1\% precision in distinguishing active signing from resting states.

\subsection{Future Directions}

\begin{enumerate}
    \item \textbf{Translation Lag Compensation}: Implementing the estimated 3.36-second décalage offset.
    \item \textbf{Skeletal Feature Extraction}: Outputting 85-point landmark sequences alongside RGB video.
    \item \textbf{SVO$\rightarrow$SOV Reordering}: Integrating syntactic reordering for alignment.
    \item \textbf{CTC Alignment}: Implementing weakly-supervised temporal alignment.
    \item \textbf{Cross-Lingual Transfer}: Leveraging bilingual (Sinhala/Tamil) nature of broadcasts.
    \item \textbf{Large-Scale Deployment}: Processing one year of broadcasts ($\sim$864,000 clips).
    \item \textbf{Baseline Models}: Training CSLR and SLT models using I3D, SlowFast, and Transformer architectures.
\end{enumerate}

\bibliographystyle{IEEEtran}
\bibliography{references}

\end{document}
